\documentclass{report}

\usepackage[utf8]{inputenc}
\usepackage[T1]{fontenc}
\usepackage[francais]{babel}
\usepackage{graphicx}
\usepackage{hyperref}
\usepackage{textcomp}
\usepackage{amsmath}
\usepackage{geometry}

\usepackage{array,multirow,makecell}
\setcounter{secnumdepth}{3}
\setcounter{tocdepth}{3}
\setcellgapes{4pt}
\makegapedcells
\newcolumntype{R}[1]{>{\raggedleft\arraybackslash }b{#1}}
\newcolumntype{L}[1]{>{\raggedright\arraybackslash }b{#1}}
\newcolumntype{C}[1]{>{\centering\arraybackslash }b{#1}}

\usepackage[pages=some]{background}
\backgroundsetup{
	scale=1,
	color=black,
	opacity=0.1,
	angle=0,
	contents={%
		\includegraphics[width=\paperwidth,height=\paperheight]{img/ulbback.jpg}
	}%
}


\title{\Huge\emph{Cooling down a coke can}\\
	\LARGE Experiment study\\
	\vspace{11pt}
	\normalsize Fluid Mechanics and Transport Processes}

\date{December 2015}

\author{Nathan DE PRYCK \and Miquel KEGELEIRS \and Mischa MASSON}



\begin{document}
	
	\begin{figure}[t]
		\includegraphics[width=15cm]{img/entete.PNG}
	\end{figure}
	
	\maketitle
	
	\renewcommand{\abstractname}{``Cooling down a coke can'' \\Experiment study by Nathan De Pryck, Miquel Kegeleirs and Mischa Masson\\ Université Libre de Bruxelles\\2015-2016.}
	
	\BgThispage
	
	\begin{abstract}
		TO BE WRITTEN
	\end{abstract}
	
	\clearpage
	
	
	\tableofcontents
	
	\chapter{Introduction}\label{intro}
	
	This report outlines the physics behind an experiment produced during the MECA-H3001 ``Fluid Mechanics and Transport Processes'' course of the ULB (Université Libre de Bruxelles) on Friday, the sixteenth of October 2015.
	
	The experiment can be found by following \hyperref{https://www.youtube.com/watch?v=MSwc_IAPh3E}{}{}{this link}\footnote{\url{https://www.youtube.com/watch?v=MSwc_IAPh3E}} and was described as following:
	
	\begin{itemize}
		\item Three coke cans were available as well as a bucket of ice and water (at 0\textdegree C) and a drill to spin the	can inside the bucket.
		\item One of the cans was used to determine the initial temperature (16\textdegree C) of the fluid (essentially water) inside the can.
		\item One can was left inside the bucket for 60 seconds and a final temperature of 11.9°C was measured as a result of the conductive heat transfer with the surrounding fluid.
		\item One can was spun inside the bucket for 60 seconds at about 1000 rounds per minute, and a final temperature of 11\textdegree C was measured as a result of the convective and conductive heat transfer.
	\end{itemize}
	
	In the first chapter, convective and conductive processes that are related to the experiment will be described. Those phenomena will then be applied to the problem via a mathematical model, including a discussion of simplifications brought to the problem in order to simplify calculations. Lastly, a comparison between the model and reality followed by a conclusion on mathematical models will be presented.
	
	\chapter{Heat transfer processes}\label{htp}
	
	They are three different ways of transferring heat: conduction, convection and radiation. Due to its nature, radiation can be neglected for this experiment and will thus not be presented here.
	
	\section{Conduction}\label{cd}
	
	Thermal conduction is a heat transfer process without macroscopic movement of matter. It is initiated by a difference of temperature between contiguous bodies (or inside a body). This difference of temperature implies a difference of internal energy : the energy is higher in the warmer area than in the cooler. By diffusion and collisions between the particles which can be molecules in a fluid or conduction electrons in a solid, particles in the warmer area transfer kinetic energy to the other particles, making them moving or vibrating faster. This creates a heat flow from the warmer area to the cooler until the system reaches thermal equilibrium. Furthermore, conduction is an irreversible process.
	
	Conduction is described by the following general equation, which is demonstrated in Professor Jean-Marie Buchlin's course\cite{Buchlin}, Chapter 13.
	
	\begin{equation}
		\frac{\delta T}{\delta t} = \bigtriangledown . (\alpha \bigtriangledown T)+ \dot{Q}_v
	\end{equation}
	
	This equation can not be used by itself because of it's nature (second degree partial derivative equation). It thus needs conditions linked to properties of the system. Those can be geometrical, physical, temporal or border conditions.
	
	Conditions used and simplifications of the general equation above will be discussed in chapter~\ref{mm}
		
	\section{Convection}\label{cv}
	
	Convection is a heat transfer in fluid. Convection occurs when some fluid is in movement. The movement lead to an advection (heat is transported by matters when it's moving).
	
	Convection is described as following :
	\begin{equation}
		Convection = Conduction + Advection
	\end{equation}
	
	Seeing this, it is easily to understand that convection is superior than conduction in fluidsin a flux situation. Flow properties have a major impact in heat transfers.
	
	As convection depends on the flow (laminar, turbulent,...), we will discuss the equation to use in the next chapter(Mathematical model,chapter~\ref{mm}).
	
	\chapter{Mathematical model}\label{mm}
	
	The idea behind mathematical models is to create a simplified version of a problem, that is accurate enough to predict the behaviour of a system, but simple enough to be resolved with few calculations.
	
	This means that some simplifications of the equations seen before can be made, using the properties of the studied system.
	
	\section{Simplifying assumptions}\label{sa}
	
	\section{Simplified model}\label{sm}
	
	\chapter[Reality and mathematical model]{Comparison between reality and the mathematical model}\label{rvmm}
	
	\chapter{Conclusion}\label{ccl}
	
	
	\bibliography{biblio}{}
	\bibliographystyle{plain}
	
\end{document}